\chapter{Introduction}
\label{chap:introduction}

\section{Motivation}

Digital Twin systems connect physical or simulated assets with cloud-based
ingestion, processing, storage, model, event, and visualization capabilities.
A layered reference architecture makes these responsibilities easier to
reason about, but it does not decide which provider services should implement
them. AWS, Azure, and Google Cloud expose different service boundaries,
behavioral guarantees, pricing units, identity mechanisms, and operational
prerequisites. A provider assignment is therefore useful only when it is both
economically explainable and functionally executable.

Multi-cloud cost optimization makes this distinction especially important. A
nominally cheap assignment may omit a mandatory delivery guarantee, bridge,
runtime artifact, or cleanup operation. Comparing such an assignment with a
complete provider bundle would confound missing functionality with a cost
advantage. Conversely, limiting every responsibility to one provider would
avoid cross-cloud complexity but leave the economic effect of layer-wise
selection unexamined.

This thesis studies the connection between these concerns. Its object is not a
general cloud-management product, but a bounded proof of concept (PoC) that can
trace one typed Digital Twin workload from functional and monetary assessment
through provider deployment, runtime verification, and destruction.

\section{Problem Statement}

Twin2MultiCloud builds on two independently developed predecessor prototypes.
The first calculated costs for a Five-layer Digital Twin model in a browser.
The second provisioned an AWS-oriented Five-layer deployment. They did not
share a machine-readable architecture contract: the optimizer's result could
not become the Deployer's input, and the Deployer could not demonstrate that
the selected responsibilities were functionally comparable across three cloud
providers.

Connecting the prototypes requires more than an API adapter. The system needs
a versioned workload, explicit provider capability mappings, attributable
pricing evidence, an exact deployment specification, cloud authority selected
for the resolved graph, durable mutation boundaries, and verification and
cleanup records. It must also distinguish deterministic offline evidence from
claims that require supervised live-cloud execution.

The implementation must remain proportionate to a Master's thesis. Features
such as production identity federation, multi-tenancy, arbitrary architecture
plugins, multiple optimization objectives, mutable infrastructure migration,
and dashboard administration would expand the software without directly
strengthening the intended research evidence. The PoC therefore fixes one
Six-layer architecture, one cost objective, one local owner profile, and one
bounded evaluation workload.

\section{Research Objective}

The objective of this thesis is to operationalize a layered, cost-aware
Digital Twin model as a reproducible proof of concept across AWS, Azure, and
Google Cloud. Twin2MultiCloud must connect typed workload intent, functional
provider mapping, cost calculation, deployment, verification, and cleanup
without treating unlike cloud products as automatically equivalent.

The thesis addresses the following research questions:

\begin{enumerate}
    \item[\textbf{RQ1}] \textit{How can a configuration-driven platform
    operationalize a layered, cost-aware Digital Twin model into reproducible
    deployments across AWS, Azure, and Google Cloud?}

    \item[\textbf{RQ2}] \textit{How can provider-specific cloud services be
    mapped to functionally complete and cost-comparable implementations of one
    shared Digital Twin architecture contract without assuming one-to-one
    service equivalence?}

    \item[\textbf{RQ3}] \textit{To what extent can layer-wise multi-cloud
    service selection reduce estimated operational cost compared with
    functionally equivalent single-cloud baselines?}

    \item[\textbf{RQ3.1}] \textit{How do single-cloud and multi-cloud
    configurations compare under identical workload and functional
    requirements?}

    \item[\textbf{RQ3.2}] \textit{How does introducing an explicit Eventing
    and Messaging Layer affect functional coverage, architecture topology, and
    total estimated cost?}
\end{enumerate}

Refactoring the predecessor projects and classifying deployment failures
remain important engineering contributions and evidence for RQ1. They are not
treated as independent research questions because the central research object
is the operationalized, functionally gated cost comparison.

\section{Contributions}

The main contributions are:

\begin{enumerate}
    \item a standalone, versioned Six-layer Eventing contract that connects
    workload intent, cost evidence, resolved architecture, deployment
    specification, and runtime verification;
    \item a functional-completeness method that compares curated AWS, Azure,
    and GCP service bundles before monetary ranking;
    \item an evidence-backed cost-only optimizer with explicit formula, unit,
    tier, transfer, bridge, and source ownership;
    \item a Terraform-based multi-cloud deployment proof of concept with
    graph-derived readiness, immutable Twin lifecycle, durable operations, and
    cleanup evidence; and
    \item a bounded evaluation design consisting of provider-local baselines,
    all six directed provider pairs, and explicit cost and cleanup guardrails.
\end{enumerate}

The architectural refactoring of the two predecessor projects, the use of
small extension patterns without inactive product features, and the unified
Management API and Flutter workflow support these contributions.

\section{Document Structure}

Chapter~\ref{chap:background} introduces Digital Twins, layered architectures,
multi-cloud computing, cost optimization, Infrastructure as Code, and the
software patterns used in the PoC. Chapter~\ref{chap:relatedWork} positions the
work against Digital Twin platforms, multi-cloud deployment systems, and cloud
cost tools. Chapter~\ref{chap:legacyAnalysis} analyzes the predecessor
prototypes and derives the research-driven transformation. Chapter~\ref{chap:systemArchitecture}
describes the final bounded architecture and its evidence chain.
Chapter~\ref{chap:evaluation} defines the deterministic and supervised
evaluation, Chapter~\ref{chap:discussion} interprets its results and validity,
and Chapter~\ref{chap:conclusion} summarizes the answers and future research.
